\documentclass[11pt,a4paper]{ctexart}
\usepackage[margin=2.2cm]{geometry}
\usepackage{xcolor}
\usepackage{enumitem}
\usepackage{amssymb}
\usepackage{booktabs}
\usepackage[normalem]{ulem}
\usepackage{array}
\usepackage{longtable}
\usepackage{titlesec}
\usepackage{hyperref}
\definecolor{navy}{HTML}{1A365D}
\hypersetup{colorlinks=true,linkcolor=navy,urlcolor=navy}
\setlist{nosep}
\titleformat{\section}{\Large\bfseries\color{navy}}{}{0pt}{}
\titleformat{\subsection}{\large\bfseries\color{navy}}{}{0pt}{}
\pagestyle{plain}
\renewcommand{\arraystretch}{1.25}

\title{\textbf{七年级下册语法总结}\\[4pt]
\large 人教版（新课纲版）七年级下册 · 按教材书末 Grammar 附录整理}
\author{初中英语资料站}
\date{2026-11-12 · 附配套专项练习 192 题}
\begin{document}
\maketitle
\section*{使用说明}
\begin{itemize}[leftmargin=1.4em]
\item 本讲义按七年级下册教材书末 Grammar 附录整理，表格与例句取自教材原文。
\item 网页版：\url{https://1901151458-prog.github.io/english-materials/g7b-grammar/}
\end{itemize}

\section{名词（Nouns）}
\noindent\textcolor{gray}{可数性、名词复数的构成与读音}\par\smallskip
\subsection*{1. 可数名词和不可数名词}

名词按其所表示的事物的性质可分为\textbf{可数名词}和\textbf{不可数名词}。可数名词有单数和复数两种形式，如 an apple–two apples、a bag–some bags。不可数名词一般只有一种形式，没有复数形式，如 milk、bread、rice 等。还有一些词既可以充当可数名词，也可以充当不可数名词，如 fruit、fish、chicken 等。

\subsection*{2. 名词复数}

名词的复数形式一般在单数形式后面加 \textbf{-s} 或 \textbf{-es}。名词复数词尾 -s 或 -es 的读音规律一般是：在清辅音后读作 \textbf{/s/}，在浊辅音和元音后读作 \textbf{/z/}，在 /s/、/z/、/ʃ/、/tʃ/、/dʒ/ 等后读作 \textbf{/ɪz/}。

{\footnotesize
\begin{longtable}{p{4.47cm}p{4.47cm}p{4.47cm}}
\hline
类别 & 构成方法 & 例词与读音 \\
\hline
\endfirsthead
\hline
类别 & 构成方法 & 例词与读音 \\
\hline
\endhead
\hline
\endfoot
一般情况 & 词尾加 \textbf{-s} & book–books /bʊks/\\ cap–caps /kæps/\\ job–jobs /dʒɒbz/\\ bag–bags /bæɡz/\\ key–keys /kiːz/\\ boy–boys /bɔɪz/ \\
以 s、x、ch、sh 等结尾的词 & 词尾加 \textbf{-es} & class–classes /ˈklɑːsɪz/\\ box–boxes /ˈbɒksɪz/\\ watch–watches /ˈwɒtʃɪz/\\ dish–dishes /ˈdɪʃɪz/ \\
以辅音字母加 y 结尾的词 & 变 y 为 i，再加 \textbf{-es} & family–families /ˈfæməlɪz/\\ strawberry–strawberries /ˈstrɔːbəriz/\\ party–parties /ˈpɑːtɪz/ \\
\end{longtable}
}

\begin{center}\fbox{\parbox{0.92\textwidth}{\small \textbf{教材注释：}少数名词的复数形式是不规则的，如 \textbf{man–men、woman–women、child–children、sheep–sheep} 等。}}\end{center}

\section{形容词（Adjectives）}
\noindent\textcolor{gray}{作定语、表语、补语}\par\smallskip
形容词可修饰名词或代词，在句子中通常用作\textbf{定语、表语和补语}等。

\subsection*{1. 用作定语}

\begin{itemize}[leftmargin=1.4em]
\item But only \underline{clever} people can see them!\\ 但只有聪明人才能看见它们！
\item It feels like a \underline{magical} place.\\ 感觉像个神奇的地方。
\item Those \underline{young} people like sport very much.\\ 那些年轻人很喜欢运动。
\end{itemize}

\subsection*{2. 用作表语（置于系动词之后）}

\begin{itemize}[leftmargin=1.4em]
\item She is \underline{busy}. The film is \underline{boring}.\\ 她很忙。电影很无聊。
\item A good student must be \underline{diligent}.\\ 好学生必须勤奋。
\item They are very \underline{hard-working}.\\ 他们很努力。
\end{itemize}

\subsection*{3. 用作补语}

\begin{itemize}[leftmargin=1.4em]
\item We need to keep our school \underline{clean and tidy}.\\ 我们要保持学校干净整洁。
\item She made me \underline{angry} last night.\\ 她昨晚让我生气了。
\end{itemize}

\begin{center}\fbox{\parbox{0.92\textwidth}{\small \textbf{三种位置一句话记：}定语在名词\textbf{前}（clever people）；表语在系动词\textbf{后}（is busy）；补语在宾语\textbf{后}（made me angry）。}}\end{center}

\section{副词（Adverbs）}
\noindent\textcolor{gray}{用法与意义、频度副词}\par\smallskip
\subsection*{1. 副词的用法和意义}

副词在句中主要用作\textbf{状语}，修饰动词、形容词、其他副词以及全句，表示时间、地点、程度、方式等概念。

{\footnotesize
\begin{longtable}{p{6.70cm}p{6.70cm}}
\hline
修饰对象 & 例句 \\
\hline
\endfirsthead
\hline
修饰对象 & 例句 \\
\hline
\endhead
\hline
\endfoot
修饰动词 & It’s raining \textbf{heavily}. \\
修饰形容词 & The dog is \textbf{quite} cute. \\
修饰其他副词 & She drives \textbf{very carefully}. \\
修饰全句 & \textbf{Usually} I have bread and milk for breakfast. \\
\end{longtable}
}

\subsection*{2. 频度副词}

频度副词是副词的一种，表示事情发生的\textbf{频率}。常用的有 \textbf{always、usually、often、sometimes、seldom (hardly ever)、never} 等，它们在频度上的差异为：

\begin{center}\fbox{\parbox{0.92\textwidth}{\small \textbf{always > usually > often > sometimes > seldom (hardly ever) > never}}}\end{center}

频度副词通常放在\textbf{实义动词前面}。如果句子里有\textbf{情态动词、助动词或系动词 be}，频度副词通常放在这些动词\textbf{之后}。

\begin{itemize}[leftmargin=1.4em]
\item I \underline{always play} tennis near my house.\\ 实义动词前。
\item We \underline{usually practise} three times a week after school.\\ 实义动词前。
\item I \underline{often go swimming} because it keeps me healthy.\\ 实义动词前。
\item A kind word is \underline{never lost}.\\ 系动词 be 之后（被动结构）。
\end{itemize}

\section{情态动词（Modal Verbs）}
\noindent\textcolor{gray}{can / could、may、must、have to}\par\smallskip
情态动词表示说话人的语气和情态，如请求、许可、意愿、可能等。在形式上，情态动词一般\textbf{没有人称和数的变化}。有的情态动词没有时态的变化，如 must。情态动词本身不能单独作谓语，\textbf{必须和动词原形连用}。常见的有 can、may、must、have to、should、would 等。情态动词的否定式一般是在其后加 not 构成，一般疑问句通常将它提到句首。

\subsection*{1. can 和 could 的用法}

{\footnotesize
\begin{longtable}{p{6.70cm}p{6.70cm}}
\hline
功能 & 例句 \\
\hline
\endfirsthead
\hline
功能 & 例句 \\
\hline
\endhead
\hline
\endfoot
表示能力 & We \textbf{can speak} a little English.\\ I \textbf{can play} badminton, but I \textbf{can’t play} volleyball.\\ He \textbf{could swim} when he was only six. \\
表示许可 & You \textbf{can play} basketball at the weekend.\\ We \textbf{can’t run} in the hallways.\\ Father said I \textbf{could swim} in the pool. \\
表示请求 & \textbf{Can you help} me with my Chinese?\\ could 也可表示请求，但比 can 的语气\textbf{更委婉}：\textbf{Could you tell} him to call me back? \\
\end{longtable}
}

\subsection*{2. may 的用法}

\begin{itemize}[leftmargin=1.4em]
\item \underline{May} I borrow your dictionary?\\ 表示许可，「可以」。
\item Class is over. You \underline{may go} now.\\ 表示许可。
\item I \underline{may be} late, so don’t wait for me.\\ 表示可能性，「也许；可能」。
\item That \underline{may be} true.\\ 表示可能性。
\end{itemize}

\begin{center}\fbox{\parbox{0.92\textwidth}{\small \textbf{may 与 can：}两者都可表示许可，但 \textbf{may 的用法比 can 更正式，语气更委婉}。}}\end{center}

\subsection*{3. must 的用法}

must 表示「\textbf{必须}」。

\begin{itemize}[leftmargin=1.4em]
\item Don’t arrive late for class. We \underline{must be} on time.\\ 必须准时。
\item You \underline{must read} a book before you watch TV.\\ 必须先读书。
\end{itemize}

\subsection*{4. have to 的用法}

have to 表示「\textbf{不得不；必须}」。

\begin{itemize}[leftmargin=1.4em]
\item We \underline{have to be} quiet in the library.\\ 在图书馆必须安静。
\item She \underline{has to do} her homework first when she gets home.\\ 她回家必须先做作业。
\end{itemize}

{\footnotesize
\begin{longtable}{p{4.47cm}p{4.47cm}p{4.47cm}}
\hline
区别 & 说明 & 例句 \\
\hline
\endfirsthead
\hline
区别 & 说明 & 例句 \\
\hline
\endhead
\hline
\endfoot
语气来源 & must 通常表示说话人的\textbf{主观看法}，语气比较强烈；have to 往往强调\textbf{客观需要}。 & We \textbf{must} hurry. I don’t want to miss the beginning of the show.\\ We \textbf{have to} hurry. The train is leaving in five minutes. \\
否定式含义 & \textbf{mustn’t} 表示「不准；禁止」，而 \textbf{don’t have to} 表示「不必」。 & You \textbf{mustn’t} talk to your mother like that.\\ You \textbf{don’t have to} come if you don’t want to. \\
人称和数 & must \textbf{没有}人称和数的变化；have to \textbf{有}人称和数的变化。 & I / We / You / They \textbf{must} do something about it.\\ I \textbf{have to} finish my work today. / She \textbf{has to} work late today. \\
\end{longtable}
}

\section{物主代词（Possessive Pronouns）}
\noindent\textcolor{gray}{名词性物主代词}\par\smallskip
物主代词是表示所有关系的代词，分为\textbf{形容词性物主代词}和\textbf{名词性物主代词}两种。七年级上册已介绍了形容词性物主代词，此处继续介绍\textbf{名词性物主代词}。

{\footnotesize
\begin{longtable}{p{3.35cm}p{3.35cm}p{3.35cm}p{3.35cm}}
\hline
物主代词 & 第一人称 & 第二人称 & 第三人称 \\
\hline
\endfirsthead
\hline
物主代词 & 第一人称 & 第二人称 & 第三人称 \\
\hline
\endhead
\hline
\endfoot
形容词性（单数） & my & your & his / her / its \\
名词性（单数） & mine & yours & his / hers / its \\
形容词性（复数） & our & your & their \\
名词性（复数） & ours & yours & theirs \\
含义 & 我的 / 我们的 & 你的 / 你们的 & 他（她、它）的 / 他们的 \\
\end{longtable}
}

\begin{center}\fbox{\parbox{0.92\textwidth}{\small \textbf{关键区别：}名词性物主代词\textbf{相当于名词}，\textbf{不能用于名词前}，说话时\textbf{要加重语气}。}}\end{center}

\begin{itemize}[leftmargin=1.4em]
\item A: Is this \underline{your} pen?\\ B: No, it isn’t. The blue pen is \underline{mine}.\\ your 后接名词；mine 独立使用。
\item A: There is a ping-pong bat here. Is it \underline{yours}?\\ B: No, it isn’t \underline{mine}. Is it Ella’s?\\ C: No, \underline{hers} is at home.\\ yours / mine / hers 都独立作成分。
\end{itemize}

\section{现在进行时（Present Continuous Tense）}
\noindent\textcolor{gray}{be + 动词-ing；动词-ing 的构成}\par\smallskip
\subsection*{1. 现在进行时的意义}

\begin{itemize}[leftmargin=1.4em]
\item 表示\textbf{现在（说话时）正在进行}的动作或存在的状态：The students are listening to the teacher. / He is sleeping now.
\item 也可表示\textbf{当前一段时间或现阶段正在进行}的动作或存在的状态：We are working on a farm these days. / I’m reading a history book this month.
\end{itemize}

\subsection*{2. 现在进行时的构成}

现在进行时是由助动词 \textbf{be（am / is / are）加动词 -ing 形式}构成的。以动词 work 为例：

{\footnotesize
\begin{longtable}{p{6.70cm}p{6.70cm}}
\hline
肯定式 & 否定式 \\
\hline
\endfirsthead
\hline
肯定式 & 否定式 \\
\hline
\endhead
\hline
\endfoot
I \textbf{am working}.\\ He / She \textbf{is working}.\\ We / You / They \textbf{are working}. & I \textbf{am not working}.\\ He / She \textbf{is not working}.\\ We / You / They \textbf{are not working}. \\
\end{longtable}
}

{\footnotesize
\begin{longtable}{p{6.70cm}p{6.70cm}}
\hline
疑问式 & 简略答语 \\
\hline
\endfirsthead
\hline
疑问式 & 简略答语 \\
\hline
\endhead
\hline
\endfoot
\textbf{Am I working?} & Yes, you are. / No, you are not (aren’t) / No, you’re not. \\
\textbf{Are you working?} & Yes, I am. / Yes, we are. / No, I am not. / No, we are not (aren’t). \\
\textbf{Is he / she working?} & Yes, he / she is. / No, he / she is not (isn’t). \\
\textbf{Are they working?} & Yes, they are. / No, they are not (aren’t). \\
\end{longtable}
}

\subsection*{3. 动词 -ing 形式的构成}

{\footnotesize
\begin{longtable}{p{4.47cm}p{4.47cm}p{4.47cm}}
\hline
类别 & 构成方法 & 例词 \\
\hline
\endfirsthead
\hline
类别 & 构成方法 & 例词 \\
\hline
\endhead
\hline
\endfoot
一般情况 & 直接加 \textbf{-ing} & go–going、ask–asking、look–looking、study–studying \\
以不发音字母 e 结尾的动词 & 去掉 e，再加 \textbf{-ing} & write–writing、make–making、take–taking \\
以重读闭音节结尾、末尾只有一个辅音字母 & 双写词尾的辅音字母，再加 \textbf{-ing} & get–getting、sit–sitting、put–putting、stop–stopping、begin–beginning \\
以 ie 结尾的动词 & 变 ie 为 y 再加 \textbf{-ing} & lie–lying、die–dying \\
\end{longtable}
}

\section{一般过去时（Simple Past Tense）}
\noindent\textcolor{gray}{意义与构成}\par\smallskip
\subsection*{1. 一般过去时的意义}

\textbf{（1）}表示过去某个时间或某一段时间内发生的动作或存在的状态，常和表示过去的时间状语连用，如 yesterday、last night、in 1990、two days ago 等。

\begin{itemize}[leftmargin=1.4em]
\item He \underline{got up} early yesterday.\\ 昨天他起得早。
\item I \underline{visited} my grandparents last week.\\ 上周我拜访了祖父母。
\end{itemize}

\textbf{（2）}也可表示过去经常或反复发生的动作，常和 often、always 等表示频度的状语连用。

He \underline{always went} to work by bus last year.\\ 去年他总是坐公交车上班。

\subsection*{2. 一般过去时的构成}

以动词 be 和 work 为例：

{\footnotesize
\begin{longtable}{p{4.47cm}p{4.47cm}p{4.47cm}}
\hline
动词 & 肯定式 & 否定式 \\
\hline
\endfirsthead
\hline
动词 & 肯定式 & 否定式 \\
\hline
\endhead
\hline
\endfoot
be & I / He / She \textbf{was} at home yesterday.\\ We / You / They \textbf{were} at home yesterday. & I / He / She \textbf{was not (wasn’t)} at home yesterday.\\ We / You / They \textbf{were not (weren’t)} at home yesterday. \\
work & I / You / He / She / We / They \textbf{worked} in a hospital last year. & I / You / He / She / We / They \textbf{did not (didn’t) work} in a hospital last year. \\
\end{longtable}
}

\begin{center}\fbox{\parbox{0.92\textwidth}{\small \textbf{两个关键：}① be 动词有 was / were 两种形式；② 实义动词的否定用 \textbf{didn’t + 动词原形}（不是 didn’t worked）。}}\end{center}

\section{配套专项练习}
\noindent\textcolor{gray}{四个专项各 48 题、共 192 题，支持在线作答与即时解析}\par\smallskip
\paragraph{名词专项 · 48 题}

\paragraph{形容词与副词专项 · 48 题}

\paragraph{物主代词与情态动词专项 · 48 题}

\paragraph{时态专项（进行时 + 过去时）· 48 题}

\begin{center}\fbox{\parbox{0.92\textwidth}{\small \textbf{建议顺序：}先按本页 01–07 讲通规则 → 再进入对应专项练习自测 → 错题回到小节复习。四个专项合计 \textbf{192 题}。\\ \textbf{考前速查：}一页纸速查表 →（A4 单页，可打印）}}\end{center}

\end{document}
